\bookchapter{Глава 22: Как AI работает в вашем телефоне}

\chapterintro{Добро пожаловать на самый крутой «магический шоу» телефона, где главная звезда — искусственный интеллект. Фотографии «щёлкнул — и готово», голосовой помощник, которому достаточно сказать пару слов, клавиатура, которая понимает вас всё лучше — за всем этим стоит невидимый мозг AI.}

\booksection{AI-камера: ваша «тёмная комната» переехала в телефон}

Современная AI-вычислительная фотография — это как если бы команда профессионального ретушёра и фотографа была встроена прямо в кнопку спуска. Прежде чем вы нажмёте кнопку, AI уже работает: сканирует сцену, распознаёт «это голубое небо», «это лицо», «это ночная съёмка».

\booksection{Голосовой помощник: как ваш «секретарь на голосе» вас понимает?}

\begin{enumerate}
  \item \textbf{Слышит (распознавание речи)}: преобразует звуковые волны в текст.
  \item \textbf{Понимает (понимание естественного языка)}: уясняет намерение, извлекает ключевые сущности.
  \item \textbf{Выполняет и отвечает}: вызывает соответствующее приложение или сервис, затем «произносит» ответ через синтез речи.
\end{enumerate}

\booksection{Локальный AI vs. облачный AI}

Локальный (на устройстве) AI выполняет вычисления на телефоне — быстро, приватно, работает офлайн. Облачный AI требует подключения к сети, но способен решать более сложные задачи.

\booksection{Большие модели AI и телефон: ваш карманный интеллектуальный помощник}

В 2025–2026 годах взрывной рост больших моделей AI (DeepSeek R1, ChatGPT-5, ERNIE 5.0, Qwen 3 и др.) глубоко изменил способ использования телефона. Телефон больше не просто инструмент выполнения команд — он стал интеллектуальным партнёром, способным понимать, рассуждать и творить. Локальные AI-чипы (NPU с производительностью свыше 80 TOPS) впервые позволили выполнять инференс больших моделей прямо на телефоне, без подключения к сети.

\booksubsection{Что такое большая модель AI?}

Большая модель AI (Large Language Model, LLM) — это AI-система, обученная на огромных объёмах текстовых данных, способная понимать и генерировать естественный язык. Представьте себе «суперстудента», прочитавшего практически весь интернет — он может отвечать на вопросы, писать статьи, переводить, программировать, анализировать данные.

\booksubsection{Применения больших моделей AI на телефоне:}

\begin{itemize}
  \item \textbf{Диалоговые помощники}: DeepSeek, ChatGPT, Doubao, Kimi и другие приложения — можно в любой момент задать вопрос, перевести, написать текст, подвести итоги статьи. Они умнее традиционных поисковиков — дают не ссылки, а готовые ответы.
  \item \textbf{Написание и творчество}: AI поможет написать письмо, подвести итоги совещания, отшлифовать текст, сгенерировать креативный копирайт.
  \item \textbf{Генерация и редактирование изображений}: AI в один клик убирает прохожих с фото, расширяет границы кадра, меняет стиль или даже генерирует картинку по текстовому описанию.
  \item \textbf{Перевод в реальном времени}: перевод во время звонка, при личной беседе, перевод по фото — за границей языковой барьер больше не проблема.
  \item \textbf{Помощь в учёбе}: AI может объяснять понятия, составлять задания, проверять домашнюю работу — это репетитор в кармане.
\end{itemize}

\booksubsection{Локальные большие модели: телефон сам запускает AI}

Новейшие флагманские смартфоны уже могут локально запускать компактные большие модели AI (обычно около 7B параметров), и вот что это меняет:

\begin{itemize}
  \item \textbf{Не нужен интернет}: AI-функции работают без сети — в самолёте, в метро, где угодно.
  \item \textbf{Выше приватность}: ваши диалоги и данные не отправляются в облако, полностью обрабатываются на телефоне.
  \item \textbf{Мгновенный отклик}: нет сетевой задержки, AI отвечает почти моментально.
  \item \textbf{Текущие ограничения}: возможности локальных моделей пока уступают облачным; сложные задачи по-прежнему требуют подключения. Но по мере роста AI-производительности чипов локальные модели будут становиться всё мощнее.
\end{itemize}

\booksubsection{Системная интеграция AI:}

\begin{itemize}
  \item \textbf{Apple Intelligence}: AI-система, появившаяся в iOS 18, интегрирует инструменты письма, «волшебную палочку» для изображений, улучшенный Siri, сводку уведомлений и др. Приоритет — локальная обработка, в облако отправляются только сложные задачи.
  \item \textbf{Pangu (Huawei)}: интегрирована в HarmonyOS, помощник Xiaoyi получил возможности большой модели, поддерживает сводку документов, «умное распознавание экрана», AI-ретушь фото и др.
  \item \textbf{Xiaomi Surge AI}: системные AI-возможности: AI-поиск по фото, AI-сводка, AI-субтитры в реальном времени, диалог с локальной большой моделью и др.
  \item \textbf{Samsung Galaxy AI}: перевод звонка в реальном времени, AI-сводка заметок, генерация и редактирование изображений, поиск обводкой и др.
\end{itemize}

\booksection{Типичные сценарии использования на практике}

Чтобы вы наглядно ощутили «суперсилу» телефонного AI, рассмотрим пять конкретных сценариев — как AI «приходит на помощь» в повседневной жизни:

\booksubsection{Сценарий 1: Перевод в реальном времени — карманный переводчик в путешествии}

\textbf{Как это работает}: В Париже не знаете французского для заказа в ресторане? Приложение «Перевод» на iPhone поддерживает режим двустороннего диалога — каждый участник видит субтитры на своём языке, вплоть до синхронного перевода. У Samsung Galaxy AI на Android есть перевод звонка в реальном времени: вы говорите по-китайски, собеседник слышит по-английски, и наоборот.

\textbf{Принцип}: распознавание речи (ASR) машинный перевод (NMT) синтез речи (TTS) — трёхступенчатый конвейер, сквозная задержка уже менее 1 секунды.

\booksubsection{Сценарий 2: AI-удаление — спасение туристического фото, «испорченного» прохожими}

\textbf{Как это работает}: Сфотографировали красивую достопримечательность, но фон забит другими туристами? Обведите их «волшебным ластиком», телефон автоматически распознает и «закрасит» прохожих, а затем интеллектуально заполнит фон — результат часто поражает.

\textbf{Принцип}: модель Inpainting сначала распознаёт и удаляет выделенную область, затем по окружающим пикселям и семантическому пониманию «домысливает» правдоподобное заполнение. Обучающая выборка таких AI-модель насчитывает миллионы изображений.

\booksubsection{Сценарий 3: AI-ретушь — один клик до «мастера фотографии»}

\textbf{Как это работает}: Получилось мрачное фото пасмурного дня? AI «магическое небо» одним кликом заменит небо на ясное, закатное или даже звёздное. AI-расширение кадра превращает вертикальную композицию в горизонтальную — телефон «домысливает» сцену за пределами кадра.

\textbf{Принцип}: основано на генеративном AI (Generative AI) и диффузионных моделях (Diffusion Model) — модель понимает семантику изображения и генерирует новые пиксели. Это текущий «передовой край» AI-обработки изображений.

\booksubsection{Сценарий 4: Обработка документов — множитель эффективности для офисных сотрудников}

\textbf{Как это работает}: После совещания загрузите аудиозапись в AI — он автоматически составит структурированные протокол, извлечёт задачи, сгенерирует сводку. 10-страничный PDF-отчёт на английском AI за 10 секунд выжимает до ключевых выводов.

\textbf{Принцип}: большая модель с длинным контекстом (Long Context LLM) способна «прочитать» десятки тысяч знаков, провести многоходовые рассуждения и извлечь информацию. Apple Intelligence, Xiaoyi (Huawei), WPS AI, Tencent Yuanbao и др. обладают этой возможностью.

\booksubsection{Сценарий 5: AI-поиск по фото и распознавание объектов — ваш «визуальный поисковик»}

\textbf{Как это работает}: Увидели на улице незнакомый цветок — навели камеру, AI сразу говорит: «Это гибискус». Ищете в галерее прошлогодние фото с моря — вводите «море», AI автоматически находит все снимки с морем, не нужно листать вручную.

\textbf{Принцип}: модель компьютерного зрения (CV) векторизует содержимое изображения (Embedding), затем сопоставляет с текстовым запросом по сходству. Такой «мультимодальный поиск» — большой скачок в управлении фотоальбомом телефона.

\booksection{Оптимизация промптов: «формула общения», чтобы AI слушался вас}

Качество ответов AI во многом зависит от того, как вы «задаёте вопрос». Один и тот же вопрос, но заданный по-разному, даёт совершенно разные ответы. Ниже мы представим\textbf{самую практичную и простую схему общения «роль–контекст–задача–ограничения»}с множеством примеров.

\begin{tcbinfo}{Зачем нужна схема?}
Многие используют AI как Baidu — вбивают несколько ключевых слов. Но AI — не поисковик, это языковая модель, которая «читает, осмысливает и отвечает».\textbf{Дайте ему больше контекста, чётко скажите, что нужно, и ограничьте от лихих уклонений}— и он выдаст точный результат. Схема переводит ваше представление о «хорошем ответе» в инструкции, понятные AI.

\end{tcbinfo}

\booksubsection{Роль (Role): скажите AI, кто он}

Дайте AI конкретную «идентичность», и он задействует профессиональные знания и стиль соответствующей области.\textbf{Чем конкретнее роль, тем профессиональнее вывод}.

\booksubsection{Практические примеры:}

\begin{itemize}
  \item «Вы —\textbf{инженер по ремонту телефонов Xiaomi с 10-летним стажем}.» (технический вопрос)
  \item «Вы —\textbf{учитель физики в средней школе, который умеет объяснять сложные вещи просто}.» (популяризация науки)
  \item «Вы —\textbf{книжный редактор с хорошим вкусом, начитанный и со своим стилем}.» (написание текстов)
  \item «Вы —\textbf{кулинарный блогер, который каждый день готовит завтраки для офисных работников}.» (повседневные вопросы)
  \item «Вы —\textbf{терпеливый и мягкий врач китайской медицины, который никогда не смеётся над новичками}.» (здоровье)
\end{itemize}

\textbf{Почему это работает?}Роль задаёт AI\textbf{область знаний},\textbf{стиль изложения}и\textbf{глубину профессионализма}. Без роли AI — «универсал, но мастер ничего»; с ролью он «входит в образ» и думает как соответствующий профессионал.

\booksubsection{Контекст (Background): расскажите AI о своей ситуации}

Заранее поясните «кто вы, какова ваша текущая ситуация, какие у вас ограничения» — и AI выдаст ответ,\textbf{действительно подходящий именно вам}, а не общие фразы.

\booksubsection{Практические примеры:}

\begin{itemize}
  \item «Мой телефон —\textbf{Xiaomi 14, полгода, при игре сильно греется}.» (ситуация с устройством)
  \item «Я —\textbf{60-летняя пенсионерка, зрение слабое, со смартфоном на «вы»}.» (портрет пользователя)
  \item «Я\textbf{готовлюсь к экзамену CET-6 по английскому в июне 2026, база средняя}.» (учебный контекст)
  \item «Я владею\textbf{маленькой кофейней в городе третьего уровня, месячная выручка 30 тыс. юаней, два сотрудника}.» (бизнес-контекст)
  \item «Я\textbf{недавно родила, дома на послеродовом, катастрофически не высыпаюсь}.» (этап жизни)
\end{itemize}

\textbf{Почему это работает?}Контекст сообщает AI,\textbf{«для кого»}и\textbf{«при каких условиях»}работать. Одну и ту же тему студенту и родителю объясняют по-разному; жителю мегаполиса и жителю посёлка предложат разные решения. Без контекста AI может дать лишь ответ «наибольшего общего делителя» — обычно ни горячо ни холодно.

\booksubsection{Задача (Task): чётко укажите AI, что делать}

Начните с\textbf{глагола}(напишите, перечислите, сравните, переведите, объясните, подведите итоги…), опишите задачу\textbf{конкретно и исполнимо}. Размытая задача — размытый ответ.

\booksubsection{Практические примеры:}

\begin{itemize}
  \item «\textbf{Перечислите 5 разрешений, которые мне следует немедленно отключить}, и объясните почему.» (действие чётко)
  \item «\textbf{Переведите следующее деловое письмо с английского на русский}, сохранив деловой тон.» (задача конкретна)
  \item «\textbf{Сравните четыре телефона в ценовой категории 3000 юаней}, представьте таблицу плюсов и минусов.» (формат ясен)
  \item «\textbf{Объясните «что такое 5G» языком, понятным младшекласснику}.» (стиль задан)
  \item «\textbf{Напишите три сценария коротких видео для моего завтрака}, отдельно для офисных работников, пар и родителей с детьми.» (сегментация ясна)
\end{itemize}

\textbf{Почему это работает?}Часть «задача» — «командный центр» действий AI.\textbf{Чем яснее глагол и конкретнее требование, тем лучше AI понимает, «что сдать»}. Размытые «сделай что-нибудь», «придумай», «посмотри» — табу для AI: он начнёт угадывать ваши намерения и часто промахнётся.

\booksubsection{Ограничения (Constraints): скажите AI, чего НЕ делать}

«Ограничения» звучит как «стеснение», но на деле это мощный инструмент\textbf{сфокусировать вывод AI на том, что подходит именно вам}. Типичные параметры: объём, формат, тон, чего избегать, что обязательно включить.

\booksubsection{Практические примеры:}

\begin{itemize}
  \item «\textbf{Не более 200 знаков, перечислением}.» (объём + формат)
  \item «\textbf{Тон дружелюбный и естественный, как в разговоре с другом}, без «во-первых, во-вторых, в-третьих» — казёнщины.» (тон)
  \item «\textbf{Без профессионального жаргона}, а если термин нужен — сразу объясните.» (ограничение лексики)
  \item «\textbf{Каждый абзац не более 3 строк}, удобно читать с телефона.» (верстка)
  \item «\textbf{Обязательно включите: сценарий использования, плюсы и минусы, соответствие моему бюджету}.» (обязательные элементы)
  \item «\textbf{Не давайте конкретных инвестиционных рекомендаций}, только объясните принцип.» (граница)
\end{itemize}

\textbf{Почему это работает?}Без ограничений AI легко «свободно импровизирует» — либо слишком академично, либо слишком длинно, либо уходит от темы. Ясные ограничения = ясные «критерии сдачи», ожидания сторон совпадают, и результат предсказуем.

\booksubsection{Оценка качества промпта: 5 признаков хорошего промпта}

Написать хороший промпт — не разовое действие.\textbf{Сначала проверьте себя по этим пяти пунктам}, потом отправляйте AI:

\begin{enumerate}
  \item \textbf{Конкретно, а не абстрактно}: не «напиши статью», а «напиши руководство по выбору телефона для молодых мам, 800 знаков, 3 раздела».
  \item \textbf{С контекстом, а не изолированно}: расскажите AI «кто я и зачем спрашиваю» — тогда ответ будет «заточен под вас».
  \item \textbf{С ясным форматом вывода}: таблица, список, абзацы или JSON? Скажите заранее.
  \item \textbf{Ограничения на месте}: объём, стиль, чего избегать — хотя бы один из трёх уже даёт заметное улучшение.
  \item \textbf{Проверяемо AI}: после ответа попросите AI пересказать ваши требования; если не может — значит, вы не были достаточно ясны.
\end{enumerate}

\booksubsection{Частые проблемы и решения}

\begin{center}
\begin{tabularx}{\linewidth}{|X|X|}
\hline
\textbf{Симптом} & \textbf{Решение} \tabularnewline
\hline
Ответ слишком общий, непрактичный & Добавьте контекст (кто вы / какой сценарий) \tabularnewline
Ответ слишком академичный / казённый & Добавьте ограничение тона + список «казённых» фраз \tabularnewline
Ответ ушёл от темы / не о том & Переформулируйте задачу через глагол и потребуйте «отвечать только на это» \tabularnewline
Ответ слишком длинный & Жёстко ограничьте объём (до 200 знаков) \tabularnewline
AI уверенно говорит ерунду & Добавьте «если не уверены — пометьте 'не уверен'» + проверка ключевых фактов \tabularnewline
Первая версия не устраивает & Не переписывайте промпт, добавьте в конец «сделай ещё …» \tabularnewline
\hline
\end{tabularx}
\end{center}

\booksubsection{Три полных сценария}

\begin{tcbinfo}{Сценарий 1: Помощник в быту (варианты завтраков)}
\textbf{Промпт}:

Вы — кулинарный блогер, специалист по питательным завтракам (\textbf{роль}). Я — работающая мама, у которой каждое утро есть лишь 15 минут на завтрак; дома 6-летний ребёнок и 60-летняя свекровь (\textbf{контекст}).\textbf{Составьте 7 неповторяющихся вариантов завтраков}(\textbf{задача}). Каждый завтрак: (1) время подготовки $\leq$ 15 мин (2) сбалансирован (белки + углеводы + овощи/фрукты) (3) подходит всем возрастам (4) продукты доступны в обычном супермаркете (\textbf{ограничения}).

\textbf{Почему эффективно}: роль задаёт профессионализм, контекст — общие черты трёх типов едоков, задача чётко определяет результат, 4 ограничения обеспечивают исполнимость.

\end{tcbinfo}

\begin{tcbinfo}{Сценарий 2: Помощь в учёбе (подготовка к CET-6)}
\textbf{Промпт}:

Вы —\textbf{опытный преподаватель CET-6 по английскому}(роль). Я — обычный бакалавр,\textbf{CET-4 — 480 баллов, цель CET-6 — 500, осталось 3 месяца}(контекст).\textbf{Составьте для меня исполнимый 12-недельный план подготовки}(задача). Требования: (1) каждую неделю — распределение времени по навыкам (аудирование, чтение, письмо, устная речь) (2) каждый месяц — ключевое направление прорыва (3) порекомендуйте 3 подходящих пособия (с обоснованием на китайском) (4) вывод в формате таблицы (ограничения).

\textbf{Почему эффективно}: роль задаёт профессиональный фон, контекст — балльные и временные ориентиры, задача чётко определяет «поэтапный план», ограничения требуют «табличного» и «исполнимого» формата.

\end{tcbinfo}

\begin{tcbinfo}{Сценарий 3: Рабочая эффективность (подведение итогов совещания)}
\textbf{Промпт}:

Вы —\textbf{опытный административный помощник}(роль). Я только что провёл 1-часовое совещание отдела,\textbf{обсуждавшее план маркетинга на Q2}(контекст).\textbf{Оформите следующую запись совещания в структурированный протокол}(задача). Требования: (1) разделите на «решения / задачи / риски и возражения» (2) для каждой задачи укажите ответственного и дедлайн (если не указано — пишите «уточнить») (3) общий объём $\leq$ 600 знаков (4) формат markdown (ограничения).

\textbf{Почему эффективно}: роль задаёт административную перспективу, контекст — тему совещания, задача чётко определяет структурированный вывод, 4 ограничения делают результат готовым к повторному использованию.

\end{tcbinfo}

\textbf{Продвинутые приёмы: сделайте AI своим «личным помощником»}\par

Попробуйте эти продвинутые подходы, чтобы ощутить настоящую силу AI:

\begin{itemize}
  \item \textbf{Многораундовая оптимизация}: если ответ AI не устраивает, скажите: «Это слишком профессионально, объясните ещё раз языком, понятным школьнику».
  \item \textbf{Цепочечное рассуждение (CoT)}: попросите AI «рассуждать шаг за шагом» — качество ответа обычно выше: «Проанализируйте эту математическую задачу по шагам».
  \item \textbf{Мульти-AI-коллаборация}: A пишет черновик, B шлифует, C переводит — комбинация работает лучше.
  \item \textbf{Системный промпт}: в настройках AI заранее укажите свою личность, предпочтения, запреты — ответы станут точнее.
\end{itemize}

\booksection{Ограничения и вызовы AI: сохраняйте ясность суждений}

AI мощен, но далеко не «всемогущ». Трезвое понимание его ограничений помогает использовать его лучше:

\begin{itemize}
  \item \textbf{Проблема «галлюцинаций»}: AI может уверенно выдумывать (технический термин — Hallucination). Он сочиняет несуществующие исторические события, ошибочные цитаты, ложную статистику. Важную информацию обязательно проверяйте вручную.
  \item \textbf{Устаревание знаний}: у большинства AI есть дата обрезки обучающих данных (например, начало 2024); они не знают о событиях после неё.
  \item \textbf{Ограниченные способности рассуждения}: сложное многоходовое рассуждение, математические доказательства, логические ловушки — AI по-прежнему часто ошибается.
  \item \textbf{Риск для приватности}: в облачных AI-сервисах ваши диалоги могут использоваться для обучения модели. Конфиденциальные данные лучше обрабатывать локальным AI.
  \item \textbf{Предвзятость и ценности}: AI может унаследовать предубеждения из обучающих данных — сохраняйте независимость мышления.
\end{itemize}

\begin{tcbinfo}{Советы по использованию AI}
\begin{itemize}
  \item Сгенерированный AI контент может быть неточным (проблема «галлюцинаций») — важную информацию обязательно проверяйте.
  \item Не вводите в AI личные конфиденциальные данные (номер удостоверения, пароли, банковские данные и т. д.).
  \item Грамотно используйте промпты: чем конкретнее и яснее описание, тем выше качество ответа AI.
  \item AI — инструмент, а не авторитет; сохраняйте независимое мышление и способность к суждению.
  \item При ключевых решениях (медицина, право, финансы) консультируйтесь с профессионалами, а не полагайтесь на AI.
\end{itemize}

\end{tcbinfo}

\begin{tcbpractice}{Попробуйте: три практических задания по AI}
\begin{enumerate}
  \item \textbf{Задание 1: Сложная инструкция}: скажите голосовому помощнику: «Завтра в 7 утра напомни мне взять зонт и перенеси сегодняшнюю встречу в 15:00 на 16:00». Проверьте, может ли он выполнить многошаговую операцию за один раз.
  \item \textbf{Задание 2: AI-ретушь}: сфотографируйте туристический объект с прохожими, используйте «волшебный ластик» или «AI-ластик», чтобы убрать прохожих. Оцените, как AI «домысливает» фон.
  \item \textbf{Задание 3: Оптимизация промпта}: сначала задайте простой промпт «напиши туристический путеводитель», затем составьте детальный промпт по CRISPE-схеме. Сравните качество двух ответов.
\end{enumerate}

\end{tcbpractice}

\begin{tcbinfo}{Пример 2: Полный процесс миграции данных при смене телефона}
\begin{itemize}
  \item \textbf{Сценарий}: iPhone 13 Xiaomi 15, смена бренда.
  \item \textbf{Шаги}: (1) Старый телефон: резервная копия iCloud / встроенный бэкап WeChat (2) Выйти из Apple ID + отключить «Найти» (3) Новый телефон: отсканировать QR-код в приложении «Xiaomi Migration» (4) Выбрать данные для переноса (5) WeChat: отдельная миграция «истории чатов» (6) Проверить ключевые данные (контакты / галерея / WeChat) (7) Старый телефон: зашифрованная очистка.
  \item \textbf{Ожидаемый результат}: нулевая потеря данных, старый телефон можно безопасно продать.
\end{itemize}

\end{tcbinfo}
